\documentclass{article}
\usepackage[utf8]{inputenc}
\usepackage[italian]{babel}
\usepackage[ruled,vlined]{algorithm2e}
\usepackage{amsmath}

\begin{document}

\section*{Introduzione}
Per risolvere alcuni problemi di \textbf{ottimizzazione} usiamo le \textbf{metodologie greedy},
delle tecniche risolutive dove ad ogni passaggio scegliamo la soluzione che, al momento, è la 
più promettente. Viene usata nei problemi dove abbiamo un numero $C$ di oggetti e dei vincoli che
devono essere soddisfatti. \\
La soluzione $S$ che vogliamo trovare è infatti un sottoinsieme degli oggetti di $C$, ma accetteremo
$S$ solo se rispetta i vincoli posti dal problema. 

\vspace{0.3cm}
% Personalizzazioni per riprodurre lo stile esatto dell'immagine
\SetKwProg{Fn}{algoritmo}{}{}%
\SetKwFunction{ParadigmaGreedy}{paradigmaGreedy}%
\SetKwFunction{Seleziona}{seleziona}%
\SetKwFunction{Ottimo}{ottimo}%
\SetKwFunction{Ammissibile}{ammissibile}%

\begin{algorithm}[H]
\Fn{\ParadigmaGreedy{\textnormal{\textit{insieme di candidati C}}} $\to$ \textnormal{\textit{soluzione}}}{

    \BlankLine

    \tcp{Variabile dove conserveremo la nostra soluzione}
    $Soluzione \gets \emptyset$\;

    \BlankLine

    \tcp{Continua finchè non troviamo la soluzione ottimale o finchè non ci sono più candidati}
    \While{$((\textbf{not } \Ottimo(Soluzione)) \textbf{ and } (C \neq \emptyset))$}{

        \BlankLine

        \tcp{Estraiamo un candidato}
        $x \gets \Seleziona(C)$\;
        $C \gets C - \{x\}$\;
        

        \BlankLine

        \tcp{Controlliamo se il canditato è una possibile soluzione ottimale}
        \If{$(\textbf{ammissibile}(S \cup \{x\}))$}{$S \gets S \cup \{x\}$\;}
    }

    \BlankLine

    \tcp{Usciti dal ciclo vuol dire che abbiamo trovato la nostra soluzione}
    \eIf{$(\textbf{ottimo}(S))$}{
        \textbf{return} $S$\;
    }{
      \tcp{Oppure abbiamo esaurito i candidati}
        \textbf{errore} non ho trovato soluzioni\;
    }
}
\end{algorithm}
\vspace{0.3cm}

Tra le tecniche greedy più famose troviamo: \textbf{bruteforce} e \textbf{problema dello zaino}

\end{document}